\scanpage{92}
II Considérons \unsure{le|les} cas $\alpha=\pm2$, donc $\unsure{\zeta|\rho}=\pm1$

On aura pour la partie homogène \inserted{$u_0$} de $u$
\[
\begin{aligned}
u_0\deleted{\unsure{\varepsilon|k}}&=\varepsilon(\mathrm{id}+v_0) &\varepsilon&=\pm1,\quad v_0^2=0\\
\text{\unsure{sc.|i.e.}}\qquad v_0&=\varepsilon u_0-\mathrm{id}
\end{aligned}
\]
De façon explicite
\[
\begin{aligned}
u(x,y)&=(1-\deleted{\unsure{x}}y+2\varepsilon\deleted{\unsure{y}}x,x)\\
u_0(x,y)&=(-y+2\varepsilon x,x)\\
\unsure{v|u}_0(x,y)&=(x-\varepsilon y,\varepsilon(x-\varepsilon y))=(x-\varepsilon y)(1,\varepsilon)\\
&&=\deleted{\unsure{\varepsilon|k}}a'(z)a
\end{aligned}
\]
où \deleted{$a'$ est la \unsure{for|f}} $a'\in\check V,\ a\in V$
\[
\left\{\begin{aligned}
a'(x,y)&=x-\varepsilon y\\
a&=(1,\varepsilon)
\end{aligned}\right.
\]
% archive sheet number: 91
